\section{The exact TPC determinant interface}
\label{sec:tpc46-interface}

This section records the coefficient packet to which the later scalar
and tensor arguments are applied.  There are three distinct objects in
play: the physical equality column at one fixed shift, the completed
alias family inherited from the determinant completion, and a family of
TPC packets reconstructed at genuinely different physical shifts.  The
first two share one source--mask coefficient field.  The third generally
does not.  Keeping this distinction explicit prevents a statement about
an auxiliary completion from being read as an averaged prime-pair
statement.

\subsection{Rows, outputs, and the physical terminal column}

Fix a nonzero physical shift \(h_0\in\Z\).  Let \(\cA_1\) and
\(\cA_2\) be the two inherited source-row systems.  A row
\(\rho\in\cA_s\) has the form
\begin{equation}
 \rho=(\ell_\rho,d_\rho),\qquad
 m_\rho=\ell_\rho d_\rho\asymp Q,\qquad
 \ell_\rho\asymp L_0,\qquad d_\rho\asymp D_0,
 \label{eq:tpc46-row}
\end{equation}
where \(\ell_\rho\) is prime, \(d_\rho\) is squarefree, and the integer
\(m_\rho\) determines \(\rho\).  The orbit set \(\cJ\) consists of
\(O(J)\) positive integers with \(j\asymp J\).  At the endpoint,
\begin{align}
 Q&=X^{267/400+o(1)},
 &J&=X^{133/400+o(1)},
 &QJ&\asymp X,
 \label{eq:tpc46-QJ-scales}\\
 D_0&=X^{10049/52500+o(1)},
 &L_0&=X^{99979/210000+o(1)}.
 \label{eq:tpc46-DL-scales}
\end{align}
The two output sets and their common Hilbert space are
\begin{equation}
 \cI_L=\{(L,\gamma,j):\gamma\in\cA_2, j\in\cJ\},
 \qquad
 \cI_R=\{(R,\alpha,j):\alpha\in\cA_1, j\in\cJ\},
 \qquad
 \Hh_{\rm out}=\ell^2(\cI_L\sqcup\cI_R).
 \label{eq:tpc46-output-space}
\end{equation}

Write the physical source coefficients as
\begin{equation}
 \gamma_\rho^{(s)}=\mu(d_\rho)\gamma_\rho^{\circ,(s)},
 \qquad s=1,2.
 \label{eq:tpc46-source-strip}
\end{equation}
Thus \(\gamma^{\circ,(s)}\) contains the source logarithm, smooth
cutoffs, and fixed residue data, but not the source M\"obius sign.  Let
\(\mathfrak A_{\alpha,\gamma}^{\rm act}(j)\) be the actual joint
source--mask multiplier and put
\begin{equation}
 \Lambda(n):=(\log n)\Phi_0(n/X).
 \label{eq:tpc46-Lambda}
\end{equation}
As usual, every terminal expression is set equal to zero outside the
declared positive support.  The physical left and right equality
columns are
\begin{align}
 [\mathscr T_{h_0}^{L}]_{\gamma,j}
 &:=-\sum_{\alpha\in\cA_1}
 \gamma_\alpha^{(1)}
 \mathfrak A_{\alpha,\gamma}^{\rm act}(j)
 \mu(n_{\alpha,j})\Lambda(n_{\alpha,j}),
 \label{eq:tpc46-physical-left}\\
 [\mathscr T_{h_0}^{R}]_{\alpha,j}
 &:=-\sum_{\gamma\in\cA_2}
 \gamma_\gamma^{(2)}
 \mathfrak A_{\alpha,\gamma}^{\rm act}(j)
 \mu(n_{\gamma,j})\Lambda(n_{\gamma,j}),
 \label{eq:tpc46-physical-right}
\end{align}
where
\begin{equation}
 n_{\rho,j}=m_\rho j+h_0.
 \label{eq:tpc46-physical-target}
\end{equation}
These formulas are the all-minus specialization of the TPC--43 prime
sign model and agree with the stripped coefficients used in TPC--45
\citep{WangTPC43,WangTPC45}.

Work on one nonempty dyadic terminal block
\begin{equation}
 N_-\le n_{\rho,j}\le N_+,
 \qquad N_-\asymp N_+\asymp X.
 \label{eq:tpc46-terminal-block}
\end{equation}
For \(R_0<N_-/\sqrt{N_+}\), define a common large-prime band
\begin{equation}
 \cP_{R_0}^{\sharp}
 :=\left\{p\in\mathbb P:\sqrt{N_+}<p\le \frac{N_-}{R_0}\right\}.
 \label{eq:tpc46-large-prime-band}
\end{equation}
The subscript \(R_0\) denotes a cofactor threshold and is unrelated to
the right output symbol in \(\cI_R\).

\begin{proposition}[Exact physical large-prime determinant sector]
\label{prop:tpc46-physical-determinant-sector}
Let \(\mathscr T_{h_0,\sharp}^{L/R}\) denote the part of
\(\mathscr T_{h_0}^{L/R}\) for which the terminal target is divisible
by a prime in \(\cP_{R_0}^{\sharp}\).  Then
\begin{align}
 [\mathscr T_{h_0,\sharp}^{L}]_{\gamma,j}
 &={}
 \sum_{\substack{p\in\cP_{R_0}^{\sharp},\ r,\\
                   \rho=(\ell,d)\in\cA_1\\
                   pr-\ell dj=h_0}}
 \gamma_\rho^{(1)}
 \mathfrak A_{\rho,\gamma}^{\rm act}(j)
 \mu(r)\Lambda(pr),
 \label{eq:tpc46-left-determinant-sector}\\
 [\mathscr T_{h_0,\sharp}^{R}]_{\alpha,j}
 &={}
 \sum_{\substack{p\in\cP_{R_0}^{\sharp},\ r,\\
                   \rho=(\ell,d)\in\cA_2\\
                   pr-\ell dj=h_0}}
 \gamma_\rho^{(2)}
 \mathfrak A_{\alpha,\rho}^{\rm act}(j)
 \mu(r)\Lambda(pr).
 \label{eq:tpc46-right-determinant-sector}
\end{align}
Every terminal target occurs at most once in either sum, and every
occurring cofactor satisfies
\begin{equation}
 R_0\le r<\sqrt{N_+}.
 \label{eq:tpc46-r-range}
\end{equation}
In particular, the three nonsmooth arithmetic roles in the collapsed
equation
\begin{equation}
 pr-mj=h_0
 \label{eq:tpc46-three-slot-equation}
\end{equation}
are the prime variable \(p\), the cofactor weight \(\mu(r)\), and the
actual source--mask coefficient on the row variable \(m\).  The orbit
variable \(j\) is smooth apart from the fixed residue selector recorded
in \cref{prop:tpc46-row-output-projective}.
\end{proposition}

\begin{proof}
If \(p>\sqrt{N_+}\) divides a target \(n\le N_+\), it is the unique
prime from \(\cP_{R_0}^{\sharp}\) dividing that target.  Write
\(n=pr\).  The upper endpoint in
\eqref{eq:tpc46-large-prime-band} gives \(r\ge R_0\), and
\(p>\sqrt{N_+}\) gives \(r<\sqrt{N_+}\).  Also \(p>r\), so
\(p\nmid r\).  Consequently
\begin{equation}
 -\mu(pr)=\mu(r),
 \label{eq:tpc46-target-sign-descent}
\end{equation}
including the case in which either side vanishes.  Substitution into
\cref{eq:tpc46-physical-left,eq:tpc46-physical-right} proves the two
formulas.  Finally \(m=\ell d\) determines its source row, so the pair
\((\ell,d)\) may be collapsed to one sequence on \(m\) without losing
any coefficient data.
\end{proof}

The sign in \eqref{eq:tpc46-target-sign-descent} is worth fixing: the
source coefficient \(\gamma^{(s)}\) already contains \(\mu(d)\).
Thus the determinant sector contains the two physical M\"obius roles
\(\mu(d)\) and \(\mu(r)\), while the large prime contributes the
minus sign that cancels the minus in the terminal atom.

\subsection{The inherited projective separation}

The actual multiplier has the fixed-shift structure
\begin{equation}
 \mathfrak A_{\alpha,\gamma}^{\rm act}(j)
 =\mathfrak m_{\rm act}(\alpha,\gamma)
   \Xi_{\alpha,\gamma}(j)
   \cW_{\alpha,\gamma}(j/J),
 \label{eq:tpc46-actual-multiplier}
\end{equation}
where, for \(\alpha=(\ell_\alpha,d_\alpha)\) and
\(\gamma=(\ell_\gamma,d_\gamma)\),
\begin{equation}
 \mathfrak m_{\rm act}(\alpha,\gamma)
 =\one_{\ell_\alpha\ne\ell_\gamma}
  \one_{|m_\alpha-m_\gamma|>\Delta_X}
 \one_{(d_\alpha,d_\gamma)\le G_X}.
 \label{eq:tpc46-static-mask}
\end{equation}
Here \(\Delta_X\) and \(G_X\) are the inherited row-separation and
gcd thresholds from TPC--36.  We write \(\mathfrak p_\infty\) for its
projective tensor norm and
\[
 \omega_*:=\max\{\omega(d):d\text{ occurs in a physical row}\}.
\]
Only the following displayed \(X^{o(1)}\) estimate for these inherited
objects is used below.
TPC--36 proves that the static mask has bounded projective norm
\begin{equation}
 \mathfrak p_\infty(\mathfrak m_{\rm act})
 \ll \log(2Q)(3+2\sqrt2)^{\omega_*}=X^{o(1)},
 \label{eq:tpc46-global-projective-cost}
\end{equation}
and that the fixed residue and smooth factors may be included without a
power loss \citep{WangTPC36}.  We use the following direct consequence.

\begin{proposition}[Actual row--output projective separation]
\label{prop:tpc46-row-output-projective}
Let \(H_0=|h_0|\).  For every fixed \(B\ge0\), the left coefficient has
an exact representation
\begin{align}
 &\gamma_\alpha^{(1)}
  \mathfrak A_{\alpha,\gamma}^{\rm act}(j)
 \notag\\
 &\quad={}
 \sum_{a\bmod H_0}\one_{j\equiv a\!\!\pmod{H_0}}
 \int_{\Omega_a}
 c_{\omega,a}(m_\alpha)
 q_{\omega,a}(\gamma)
 \phi_{\omega,a}(j/J)\,d\nu_a(\omega),
 \label{eq:tpc46-global-row-output-separation}
\end{align}
where \(c_{\omega,a}\) is supported on the physical row integers and
\begin{equation}
 \sum_{a\bmod H_0}\int_{\Omega_a}
 \|c_{\omega,a}\|_\infty
 \|q_{\omega,a}\|_\infty
 \|\phi_{\omega,a}\|_{C^B}
 \,d|\nu_a|(\omega)
 \ll_{B,h_0}X^{o(1)}.
 \label{eq:tpc46-global-row-output-mass}
\end{equation}
The right coefficient has the symmetric representation.  If the source
integer must instead be opened as \(m=\ell d\), then for every fixed
opposite row one has an exact expansion
\begin{equation}
 \gamma_{(\ell,d)}^{(1)}
 \mathfrak A_{(\ell,d),\gamma}^{\rm act}(j)
 =\sum_{a\bmod H_0}\one_{j\equiv a\!\!\pmod{H_0}}
  \int\sum_{\varrho}
  b_{\varrho}(\ell)f_{\varrho}(d)
  \eta_{\varrho}(j/J)\,d\nu,
 \label{eq:tpc46-fixed-column-separation}
\end{equation}
whose corresponding weighted signed mass is \(X^{o(1)}\).  The static
Ferrers part of this latter cost is \(O(1+\log(2D_0))\), although the
formal number of layers can be \(O(D_0)\).
\end{proposition}

\begin{proof}
TPC--36 gives a global separation of
\(\mathfrak A^{\rm act}\) into a source-row factor, an opposite-row
factor, and a smooth orbit factor with the cost in
\eqref{eq:tpc46-global-projective-cost}.  Multiply the source-row
factor by \(\gamma_\alpha^{(1)}\).  Its pointwise size is
\(X^{o(1)}\), because the only growing elementary factor is a source
logarithm.  Source separation permits the resulting function of
\(\alpha\) to be written as a function of \(m_\alpha\).  This proves
\cref{eq:tpc46-global-row-output-separation,eq:tpc46-global-row-output-mass}.
The last statement is the fixed-column Ferrers separation of TPC--36,
including its physical row and orbit weights.
\end{proof}

Because \(\Phi_0\) is a fixed smooth terminal cutoff, Mellin inversion
also separates the target product.  We use the convention
\begin{equation}
 \widetilde\Phi_0(t)
 :=\int_0^\infty\Phi_0(x)x^{-it}\,\frac{dx}{x},
 \qquad
 \Phi_0(x)=\frac1{2\pi}\int_{\R}
 \widetilde\Phi_0(t)x^{it}\,dt.
 \label{eq:tpc46-Mellin-convention}
\end{equation}
If
\begin{equation}
 d\sigma(t)=\frac1{2\pi}\widetilde\Phi_0(t)X^{-it}\,dt,
 \label{eq:tpc46-Mellin-measure}
\end{equation}
then
\begin{equation}
 \Lambda(pr)
 =\int_{\R}
 \left((\log p)p^{it}r^{it}
       +p^{it}(\log r)r^{it}\right)d\sigma(t).
 \label{eq:tpc46-target-Mellin-separation}
\end{equation}
Every polynomially weighted variation of \(\sigma\) is finite.  The
two logarithms cost \(O(\log X)=X^{o(1)}\) on the terminal support.

\begin{corollary}[Exact product-ready determinant layers]
\label{cor:tpc46-product-ready-layers}
Each coordinate in
\cref{eq:tpc46-left-determinant-sector,eq:tpc46-right-determinant-sector}
 is an \(X^{o(1)}\)-projective superposition of layers of the form
\begin{equation}
 q_\omega(\gamma)
 \one_{j\equiv a\!\!\pmod{H_0}}
 \sum_{\substack{p,r,m\\pr-mj=h_0}}
 a_\omega(p)b_\omega(r)c_\omega(m)
 \phi_\omega(j/J),
 \label{eq:tpc46-three-coefficient-layer}
\end{equation}
where the prime restriction is in \(a_\omega\), the factor \(\mu(r)\)
is in \(b_\omega\), and the literal source and mask coefficient is in
\(c_\omega\).  Equivalently, after opening \(m=\ell d\), the layers
have the form
\begin{equation}
 \one_{j\equiv a\!\!\pmod{H_0}}
 \sum_{\substack{p,r,\ell,d\\pr-dj\ell=h_0}}
 a_\omega(p)b_\omega(r)
 c_\omega(\ell)f_\omega(d)\phi_\omega(j/J).
 \label{eq:tpc46-opened-determinant-layer}
\end{equation}
No arithmetic coefficient has been replaced by a smooth cutoff in
either identity.
\end{corollary}

The residue selector is fixed because \(h_0\) is fixed, but it is not a
smooth function of \(j/J\).  It has the exact finite Fourier expansion
\begin{equation}
 \one_{j\equiv a\!\!\pmod{H_0}}
 =\frac1{H_0}\sum_{\tau=0}^{H_0-1}
 \e\!\left(\frac{\tau(j-a)}{H_0}\right).
 \label{eq:tpc46-residue-selector-Fourier}
\end{equation}
Thus every product-ready layer is a finite, \(O_{h_0}(1)\)-mass
superposition of fixed modulations of a smooth orbit profile.  It is
\emph{progression-smooth}, not literally an unmodulated smooth
sequence.  The progression form of the shift-spectral theorem is
proved in \cref{cor:tpc46-progression-shift-gap}: for a band-limited
profile its exponent range is unchanged because \(H_0\) is fixed.
The selector alone does not make the actual profile band-limited, and
it does not regularize the pulled-back dual field below.

This corollary is an interface theorem, not a scalar fixed-shift
estimate.  In particular, the coherent-complementary-factor theorem in
TPC--36 leaves a character-level quantity of the form
\begin{equation}
 \left|\sum_r w(r)\widehat R_r(\chi)\right|^2
 \label{eq:tpc46-inherited-coherent-gate}
\end{equation}
and proves that Cauchy--Schwarz in the coherent variable can be sharp.
Thus product-ready separation does not itself close even the scalar
fixed-\(h_0\) determinant sum.

\subsection{Residual duality and the tensor coefficient}

At the physical equality alias, the terminal support has
\((d_\rho,n_{\rho,j})=1\).  If \(n_{\rho,j}=pr\) with
\(p>\sqrt{N_+}\), then \(d_\rho,p,r\) are pairwise coprime and the
large-prime deletion leaves the residual label
\begin{equation}
 s=d_\rho r.
 \label{eq:tpc46-physical-residual-label}
\end{equation}
After deleting zero atoms, every occurring \(r\) is squarefree; since
\((d_\rho,r)=1\) and \(d_\rho\) is squarefree, every occurring
residual label \(s=d_\rho r\) is squarefree.  Let
\(\Hh_{\rm res}\) be the residual Hilbert space with coordinates
\((i,s)\), where \(i\in\cI_L\sqcup\cI_R\) and \(s\) ranges over these
supported squarefree labels.  In the notation of
TPC--45, define \(x_p\in\Hh_{\rm res}\) on the present large-prime
sector by the same coordinate rule: its \((i,s)\)-entry is the stripped
atom with full label \(sp\), and is zero when that atom is absent.  This
definition does not require the more restrictive very-high cutoff used
for the singleton-incidence theorem there.  Let
\begin{equation}
 (Uy)_{i,s}=\mu(s)y_{i,s}.
 \label{eq:tpc46-residual-unitary}
\end{equation}
Thus \(|\mu(s)|=1\) on the coordinate set, so \(U\) is unitary, and
\(\|\sum_p x_p\|=\|\sum_pUx_p\|\).

\begin{proposition}[Exact pulled-back dual tensor]
\label{prop:tpc46-dual-tensor}
For \(z=(z^L,z^R)\in\Hh_{\rm res}\), the left part of the dual pairing
is
\begin{align}
 \left\langle\sum_{p\in\cP_{R_0}^{\sharp}}Ux_p,z\right\rangle_L
 ={}-&\sum_{\substack{\gamma,j,\rho=(\ell,d)\in\cA_1,\\
                       p\in\cP_{R_0}^{\sharp},r\\
                       pr-dj\ell=h_0}}
 \gamma_\rho^{(1)}
 \mathfrak A_{\rho,\gamma}^{\rm act}(j)
 \Lambda(pr)\mu(r)
 \overline{z^L_{\gamma,j,dr}}.
 \label{eq:tpc46-left-dual-tensor}
\end{align}
The right part is
\begin{align}
 \left\langle\sum_{p\in\cP_{R_0}^{\sharp}}Ux_p,z\right\rangle_R
 ={}-&\sum_{\substack{\alpha,j,\rho=(\ell,d)\in\cA_2,\\
                       p\in\cP_{R_0}^{\sharp},r\\
                       pr-dj\ell=h_0}}
 \gamma_\rho^{(2)}
 \mathfrak A_{\alpha,\rho}^{\rm act}(j)
 \Lambda(pr)\mu(r)
 \overline{z^R_{\alpha,j,dr}}.
 \label{eq:tpc46-right-dual-tensor}
\end{align}
Consequently, after the projective separation in
\cref{cor:tpc46-product-ready-layers}, the new multiplier is precisely
\begin{equation}
 q_\omega(\gamma)\phi_\omega(j/J)
 \one_{j\equiv a\!\!\pmod{H_0}}
 \overline{z^L_{\gamma,j,d(m)r}}
 \label{eq:tpc46-pulled-back-z}
\end{equation}
on the left, with the symmetric right-hand expression.  It is a joint
coefficient in the output row, \(j\), the source divisor \(d(m)\), and
the cofactor \(r\).
\end{proposition}

\begin{proof}
For a left atom with \(n=pr\), TPC--45 gives the stripped coefficient
\begin{equation}
 B^L_{\rho;\gamma,j}
 =-\gamma_\rho^{\circ,(1)}
  \mathfrak A_{\rho,\gamma}^{\rm act}(j)\Lambda(pr).
 \label{eq:tpc46-stripped-B}
\end{equation}
The residual coordinate is \(s=dr\).  Multiplication by \(U\) inserts
\(\mu(s)=\mu(d)\mu(r)\).  Combining \(\mu(d)\) with
\(\gamma^{\circ,(1)}\) gives \(\gamma^{(1)}\), and pairing with \(z\)
gives \eqref{eq:tpc46-left-dual-tensor}.  The right calculation is
identical after exchanging the two row systems.  Substitution of
\eqref{eq:tpc46-global-row-output-separation} yields
\eqref{eq:tpc46-pulled-back-z}.
\end{proof}

The distinction from the scalar layer
\eqref{eq:tpc46-three-coefficient-layer} is exact.  An arbitrary unit
vector \(z\) need not factor in \((\gamma,j,dr)\).  Therefore
\eqref{eq:tpc46-pulled-back-z} cannot be absorbed into the \(r\)-slot,
the \(m\)-slot, or the smooth \(j\)-slot using the projective theorem
of TPC--36.  Expanding \(z\) into coordinate delta functions gives only
a crude illustration of the possible cost.  For one fixed opposite row
there are at most
\begin{equation}
 X^{o(1)}QJ=X^{1+o(1)}
 \label{eq:tpc46-fixed-output-active-count}
\end{equation}
active \((j,s)\)-coordinates, so this expansion may cost
\(X^{1/2+o(1)}\|z\|_2\).  Globally the number of atomic coordinates is
at most
\begin{equation}
 X^{o(1)}Q^2J=X^{667/400+o(1)},
 \label{eq:tpc46-global-active-count}
\end{equation}
whose square-root scale is \(X^{667/800+o(1)}\).  These are crude upper
bounds for one decomposition, not lower bounds or impossibility
theorems.  They show only that no \(X^{o(1)}\) tensor lift follows from
the inherited projective norm and \(\|z\|_2=1\).  The desired allowance
is merely \(K_B=X^{1/400+o(1)}\).

\subsection{Three shift families}

We finish by fixing terminology used later.

\begin{definition}[Inherited completed-alias family]
\label{def:tpc46-completed-alias-family}
Let \(M_{\rm al}\) be the triple-prime completion modulus of TPC--42 and put
\begin{equation}
 c_k=h_0+kM_{\rm al},\qquad k\in\cI_B.
 \label{eq:tpc46-completed-intercepts}
\end{equation}
The completed-alias column replaces \(n_{\rho,j}\) by
\(n_{\rho,j,k}=m_\rho j+c_k\), with zero extension off the terminal
support, while retaining the same row systems and the same coefficient
\(\gamma_\rho^{(s)}\mathfrak A^{\rm act}(j)\) for every \(k\).
\end{definition}

The coefficient independence in this definition is an inherited fact,
not an assumption \citep{WangTPC42,WangTPC43}.  However, for \(k\ne0\)
one may have
\begin{equation}
 g_{\rho,j,k}=(d_\rho,n_{\rho,j,k})=(d_\rho,c_k)>1.
 \label{eq:tpc46-alias-gcd}
\end{equation}
If \(n_{\rho,j,k}=pr\) is squarefree and \(p>d_\rho\), its fused Walsh
label is
\begin{equation}
 p\,\sfk(d_\rho r),
 \qquad
 \sfk(d_\rho r)=\frac{d_\rho r}{(d_\rho,r)^2},
 \label{eq:tpc46-alias-residual-label}
\end{equation}
not necessarily \(p d_\rho r\).  Accordingly, the alias-averaged dual
coefficient is
\begin{equation}
 z^L_{k,\gamma,j,\sfk(d r)}
 \label{eq:tpc46-alias-dual-coordinate}
\end{equation}
rather than the equality-alias coefficient
\(z^L_{\gamma,j,dr}\).

\begin{definition}[Frozen completed family]
\label{def:tpc46-frozen-family}
For an auxiliary integer \(c\), the frozen completed column keeps the
physical \(h_0\) row systems and source--mask coefficient field, but
replaces the determinant equation by
\begin{equation}
 pr-mj=c.
 \label{eq:tpc46-frozen-equation}
\end{equation}
All terminal factors are evaluated at \(n=pr=mj+c\) and are zero off
support.
\end{definition}

The inherited aliases \(c=c_k\) are instances of this frozen family.
For a general \(c\), it is an algebraic coefficient model.  It is not
the TPC packet that would be obtained by restarting the construction
with physical shift \(c\).

\begin{definition}[Actual local-shift family]
\label{def:tpc46-actual-local-family}
For each nonzero shift \(\nu\), the actual local-shift packet is the
packet reconstructed with physical shift \(\nu\).  Its eligible rows,
fixed residue factor \(\Xi^{(\nu)}\), source residue data, terminal
support, and joint multiplier are allowed to depend on \(\nu\).
\end{definition}

TPC--36 proves \(X^{o(1)}\) projective separation for each fixed
\(\nu\), with a constant depending on \(\nu\).  It does not prove a
uniform bound over a growing local-shift window.  Thus later spectral
identities for the frozen family become statements about actual shifted
TPC packets only after a separate uniform-provenance theorem.  No such
theorem is assumed in this paper.
